\chapter{Metodologia}

\section{Desenho da pesquisa}

O estudo foi conduzido como prova de conceito experimental de um pipeline híbrido para classificação e explicação de risco. O conjunto IEEE-CIS Fraud Detection foi usado para desenvolvimento e avaliação técnica \cite{ieee_cis_2019}. Como os dados representam comércio eletrônico/cartão e contêm atributos anonimizados, o experimento não permite estimar desempenho operacional em transações Pix reais. A transferência ao contexto Pix foi limitada a papéis analíticos explicitamente documentados.

O fluxo possui três camadas: classificação por aprendizado de máquina, explicação local por SHAP e contextualização documental por RAG. Cada camada foi definida com entrada, saída, versão e protocolo de avaliação próprios. Essa separação evita que uma narrativa produzida pelo modelo de linguagem seja confundida com a previsão do classificador ou com o conteúdo de uma norma.

\section{Preparação dos dados e protocolo experimental}

As tabelas de transação e identidade foram unidas por \texttt{TransactionID} por meio de junção à esquerda, preservando todas as transações. A análise local encontrou 590.540 transações, das quais 20.663 apresentam \texttt{isFraud=1}, correspondendo a 3,499\%. A tabela de identidade possui 144.233 linhas e cobre 24,424\% das transações do treino.

As observações foram ordenadas por \texttt{TransactionDT}, que representa um delta temporal em segundos a partir de uma referência não publicada. O valor preserva ordem e intervalos relativos, mas não foi convertido em data civil ou horário local. A avaliação principal adotou uma divisão temporal de 70\% para treino, 15\% para validação e 15\% para teste, com as transações mais antigas no treino e as mais recentes no teste. Transações com o mesmo valor de \texttt{TransactionDT} permaneceram na mesma partição. Transformações, imputação, \emph{encoding}, escalonamento, seleção e estatísticas históricas foram ajustados somente no treino.

No conjunto completo, o corte resultou em 413.378 transações de treino, 88.581 de validação e 88.581 de teste, contendo, respectivamente, 14.538, 3.042 e 3.083 registros marcados como fraude. As taxas correspondentes foram 3,517\%, 3,434\% e 3,480\%, uma diferença máxima de 0,083 ponto percentual. Portanto, a classe positiva permaneceu rara, mas sua proporção foi semelhante entre as três partições sem que o alvo fosse usado para estratificação. Ocorrências repetidas de \texttt{TransactionDT}, desconsiderando a primeira de cada valor, corresponderam a 2,9\% das linhas; ao contar todas as linhas pertencentes a grupos de \emph{timestamps} repetidos, a proporção foi 5,7\%. Nenhuma ocorrência empatada atravessou as fronteiras do corte. Uma análise aleatória estratificada pode ser apresentada apenas como complemento, sem substituir a avaliação temporal principal.

O registro \texttt{config/pix\_feature\_registry.json} documenta cada atributo derivado: fórmula, colunas-fonte, população, janela, tratamento de ausentes, controle de vazamento, tags semânticas e advertência. Em decisão registrada pela dupla em 16/08/2026, foram aprovados e implementados quatro conceitos: desvio robusto do valor no histórico proxy, frequência recente no identificador proxy, raridade de dispositivo nesse histórico e posição cíclica temporal relativa. \texttt{card1} é tratado como proxy mascarado do domínio original, nunca como conta ou chave Pix; a origem e o fuso de \texttt{TransactionDT} não são presumidos.

A regressão logística foi usada como linha de base. Ponderação de classe e SMOTE foram comparados em configurações separadas, com a reamostragem restrita à partição de treino \cite{chawla2002smote}. Para o artefato XGBoost final, o limiar foi selecionado na validação pela maximização do F1 e congelado antes da avaliação do teste. Foram registradas AUC-PR, AUC-ROC, precisão, revocação, F1 e matriz de confusão no teste preservado \cite{saito2015precision}.

\section{Camada 1: classificação de risco}

A primeira camada recebe os atributos preparados e produz a probabilidade e a classe de risco. XGBoost foi adotado como modelo principal, enquanto Random Forest e regressão logística foram usados como comparadores \cite{breiman2001random,chen2016xgboost}. O artefato XGBoost é distribuído com o pré-processador e o manifesto, e sua política de decisão foi persistida separadamente. O retorno atual do pipeline inclui classe, probabilidade e limiar, mas não expõe no mesmo pacote a versão do modelo e a versão do registro de atributos.

\section{Camada 2: explicação local por SHAP}

A segunda camada calculou contribuições SHAP para o XGBoost em transações selecionadas \cite{lundberg2017unified}. O método foi escolhido por fornecer atribuições locais aditivas e uma agregação global pela média dos valores absolutos. Depois de congelados o conjunto de dados, o pré-processamento, a ordem das colunas e o modelo, o \texttt{TreeExplainer} foi configurado com \texttt{feature\_perturbation=interventional}, fundo de 500 linhas amostradas somente do treino por \emph{seed} registrada e \texttt{model\_output=probability}. Essa escala permite verificar a soma das contribuições contra a probabilidade prevista; a saída bruta padrão do XGBoost estaria em margem ou \emph{log-odds}. A versão do SHAP foi registrada explicitamente. O Random Forest não recebeu, nesta etapa, uma análise SHAP equivalente e não foi incluído na comparação de explicações locais.

Para cada observação, foi verificada a propriedade aditiva
\begin{equation}
    f(x) \approx E[f(X)] + \sum_{j=1}^{p} \phi_j,
\end{equation}
em que $E[f(X)]$ é o valor-base e $\phi_j$ representa a contribuição do atributo $j$ na mesma escala de $f(x)$. O erro máximo de reconstrução registrado foi aproximadamente $3{,}04 \times 10^{-7}$, dentro da tolerância numérica adotada. Valores não finitos ou desalinhamento entre nomes e posições de colunas invalidam o artefato explicativo.

A análise global ordenou os atributos pela média de $|\phi_j|$ em 1.000 observações da validação. A análise local registrou casos representativos de verdadeiro positivo, falso positivo, falso negativo e verdadeiro negativo segundo o limiar histórico de 0,5. A seleção desses casos não foi refeita com o limiar final de 0,6834 e, por isso, seus rótulos locais não devem ser reinterpretados no novo ponto de operação. As contribuições são tratadas como influência estatística no comportamento do modelo, não como causalidade, culpa ou intenção.

Para conectar SHAP ao RAG, o pipeline seleciona os três fatores de maior contribuição absoluta, preservando nome, valor transformado, contribuição e direção. O desenho metodológico previa que somente atributos com semântica aprovada gerassem \emph{tags} conceituais. A implementação atual, contudo, constrói a consulta a partir dos nomes dos fatores, inclusive nomes anônimos, sem lhes atribuir significado Pix. Essa diferença foi mantida explícita para que uma limitação do protótipo não seja descrita como salvaguarda já implementada.

A avaliação verificou fidelidade numérica e alinhamento entre os atributos e suas posições. A consistência entre os valores SHAP e a narrativa posterior permanece sem resultado, pois o cliente LLM não foi implementado. A distribuição das contribuições pode depender da configuração do explicador, do conjunto de referência e da correlação entre atributos. Uma explicação fiel ao modelo não garante que o modelo esteja correto, calibrado ou livre de viés; por isso, SHAP é tratado como instrumento de auditoria do classificador e não como validação independente da decisão.

\section{Camada 3 --- Retrieval-Augmented Generation (RAG)}

A terceira camada prepara a contextualização documental da saída numérica das etapas anteriores. A arquitetura de \emph{Retrieval-Augmented Generation} combina a recuperação de trechos externos com a geração de linguagem \cite{lewis2020rag}. Neste trabalho, o RAG não altera a classe ou a probabilidade calculada pelo modelo, não substitui a explicação SHAP e não decide se houve fraude. Seu papel é recuperar contexto rastreável para a predição e para as contribuições locais, mantendo separadas a evidência do classificador, a informação documental e as limitações da análise.

A base de conhecimento é composta por quatro fontes catalogadas: o Regulamento do Pix estabelecido pela Resolução BCB nº 1/2020, a Resolução BCB nº 103/2021, o Guia do Mecanismo Especial de Devolução na versão 4.4 e a Pesquisa Febraban de Tecnologia Bancária 2024 --- volume 1 \cite{bcb_resolucao1_2020,bcb_resolucao103_2021,bcb_guia_med_v44,febraban2024}. As três primeiras fontes oferecem conteúdo normativo ou operacional, enquanto o relatório da Febraban é empregado apenas como contexto setorial. Para manter a rastreabilidade, cada documento conserva identificador, título, emissor, tipo, URL, data de publicação, versão, situação normativa, datas de vigência e hash SHA-256. Cada trecho derivado também preserva sua página ou seção e a referência ao documento de origem.

Na etapa de indexação, os arquivos PDF e JSON são extraídos e normalizados. O conteúdo é inicialmente segmentado em \emph{chunks} pais de 500 tokens, com sobreposição de 50 tokens. Como o modelo de \emph{embeddings} possui limite de 128 tokens, cada \emph{chunk} pai é dividido em janelas de até 126 tokens de conteúdo, com sobreposição de 24 tokens. As representações vetoriais são produzidas pelo modelo multilíngue \texttt{sentence-transformers/paraphrase-multilingual-MiniLM-L12-v2}, na revisão registrada pelo manifesto do índice, com vetores de 384 dimensões \cite{reimers2019sentencebert}. Após normalização L2, os vetores são armazenados em um índice FAISS exato do tipo \texttt{IndexFlatIP}; nessa configuração, o produto interno entre vetores normalizados equivale à similaridade de cosseno \cite{johnson2021billion}. A reconstrução versionada de 20 de setembro de 2026 reuniu 106 unidades extraídas, 290 \emph{chunks} pais e 1.195 janelas vetoriais. O índice, os metadados e o manifesto de integridade foram versionados. Três fontes mantiveram os mesmos hashes do registro histórico; o arquivo consolidado do Regulamento do Pix apresentou hash distinto e, por isso, a reconstrução não foi tratada como idêntica à base anterior.

A recuperação recebe uma consulta construída com a situação prevista e os nomes dos três fatores de maior contribuição absoluta identificados pelo SHAP. Nomes anônimos são preservados literalmente, sem tradução para conceitos de negócio. A consulta é codificada pelo mesmo modelo de \emph{embeddings} e comparada aos vetores do índice. O recuperador devolve os cinco trechos de maior similaridade, acompanhados de escore e metadados. A implementação não aplica limiar mínimo de similaridade nem filtro automático de vigência. Consequentemente, a presença de um trecho entre os cinco primeiros resultados não comprova sua relevância ou sua aplicabilidade temporal, e fontes com vigência escalonada exigem conferência humana antes de fundamentar uma afirmação.

A integração entre SHAP e RAG ocorre no orquestrador \texttt{src.pipeline.explicar\_transacao}. A classificação e a probabilidade são calculadas pelo XGBoost, e os três fatores de maior contribuição absoluta são encaminhados ao módulo de recuperação. O contexto do gerador preserva o nome, o sinal e o valor da contribuição SHAP, mas a versão atual do prompt não inclui o valor bruto observado da variável, a versão do modelo ou a vigência validada de cada documento. A integração, portanto, não autoriza causalidade e mantém a diferença entre o domínio IEEE-CIS e transações Pix reais, porém ainda necessita desses campos para cumprir integralmente o pacote de evidências projetado.

O prompt é montado por \texttt{src.rag.explainer.montar\_prompt}. Ele contém instruções de comportamento, a decisão e a probabilidade do classificador, os três fatores SHAP com a direção e a contribuição numérica e até cinco trechos recuperados, identificados pela fonte. As instruções exigem resposta em português do Brasil, com três a cinco frases, proíbem afirmações causais e a atribuição de significado às colunas anônimas, restringem a resposta às evidências fornecidas e determinam abstenção quando o suporte for insuficiente. Os metadados completos de URL, localização, versão, vigência e escore são mantidos no retorno estruturado do pipeline e podem ser exibidos na interface, mas não são todos incorporados ao texto atual do prompt.

O contrato do cliente de modelo de linguagem foi definido para receber o prompt e devolver texto, mas sua implementação permanece \texttt{\{\{PREENCHER\}\}}. A decisão registrada indica Claude Haiku 4.5 como opção principal e Ollama local como alternativa; nenhuma das duas integrações foi implementada ou avaliada. Na ausência do cliente, o pipeline não simula uma resposta: informa que a explicação em linguagem natural está indisponível, preserva o prompt montado para inspeção e mantém visíveis a predição, os fatores SHAP e os documentos recuperados.

As estratégias implementadas para reduzir alucinações são a fundamentação estrita, o limite do contexto aos trechos recuperados, a preservação das contribuições numéricas, a proibição de semântica inventada para variáveis anônimas e a regra de abstenção. O prompt também orienta que a transação pertence a um conjunto de pesquisa e que ninguém seja tratado como culpado. O controle automático de vigência e a validação factual pós-geração não foram implementados. A avaliação proposta considera relevância da recuperação, coerência com SHAP, fidelidade aos dados, ausência de alucinações e clareza em português \cite{es2024ragas}; entretanto, até o fechamento desta versão, somente consultas de fumaça haviam verificado o carregamento e o funcionamento do índice, sem produzir métricas de qualidade ou respostas de LLM.

\section{Avaliação integrada}

A classificação foi avaliada por métricas adequadas à classe rara, e a explicação SHAP foi verificada quanto à aditividade e ao alinhamento dos atributos. Para a recuperação, foi definido um protocolo com relevância dos trechos, precisão em $k$, \emph{recall} em $k$ e posição do primeiro resultado relevante. Para a geração, foram definidos critérios de suporte documental, coerência com SHAP, fidelidade à transação, clareza e ausência de linguagem acusatória. Como o conjunto anotado e o cliente LLM não estavam disponíveis no fechamento desta versão, a avaliação experimental de recuperação e geração permanece \texttt{\{\{RESULTADO\}\}}; consultas de fumaça não foram interpretadas como medida de qualidade. Avaliações automáticas podem apoiar, mas não substituir, a inspeção humana \cite{es2024ragas}.

\section{Reprodutibilidade, segurança e limitações}

O código, os parâmetros, as \emph{seeds}, as versões e as decisões metodológicas foram registrados no repositório. Os arquivos brutos do IEEE-CIS, a matriz de fundo do SHAP, os exemplos extraídos do conjunto e eventuais chaves de API permanecem fora do Git. Exemplos enviados a serviços externos devem ser sintéticos ou minimizados. Saídas do modelo de linguagem não devem executar comandos nem produzir decisão automática de bloqueio, acusação ou devolução.

As principais ameaças à validade são a diferença de domínio entre IEEE-CIS e Pix, a anonimização de atributos, a mudança de distribuição, a confirmação imperfeita do rótulo e a ausência de validação em uma instituição financeira. Somam-se a elas a falta de avaliação anotada do RAG, a ausência do cliente LLM e a diferença entre o limiar de 0,6834 selecionado na validação do artefato final e o valor padrão de 0,5 atualmente utilizado pela interface. Portanto, os resultados constituem evidência parcial de viabilidade técnica da integração e não caracterizam um sistema pronto para produção.
